\section{One managed station attempt}
\label{sec:managed-bw500}

The Linkwise link coordinator owns the station profile, candidate eligibility,
lease, and attempt record. The model proposes a candidate through the local
Model Context Protocol (MCP) service. Deterministic station code decides whether
to admit it. With an observation profile the service imports and ranks
evidence. With a shadow profile it also prepares a simulated radio and runs
either a scripted outcome or the managed virtual BW500 fixture. The fixture
uses two real waveform modem processes and two Winlink B2F peer clients. The
radio has zero power, simulated CAT readback, and no physical keying route.
\repolink{rust/docs/mcp-usage.md}{Station profiles and tools.}

\stationfigure[0.85\textheight]{Linkwise_Managed_Attempt}{One installed shadow attempt from ranked candidate through local evidence and cleanup.}{fig:managed-attempt}

\paragraph{Prepare, authorize, start.}
\path{link_rank_gateways} records a ranking from configured RMS inventory and
optional VOACAP forecast files. The selected ranking and channel identify a
candidate. \path{link_prepare_candidate} revalidates its eligibility, acquires
one expiring station lease, and verifies simulated dial frequency, USB mode,
zero RF power, and PTT down. The prepared handle binds the candidate and
readback revision. \path{link_request_session_authorization} refreshes the
evidence and verifies that preparation again before issuing a short-lived,
single-use grant. \path{link_start_session} checks the prepared handle,
grant, and readback before journaling one active attempt and returning its ID.
The station then runs the attempt asynchronously within the configured
session time. Repeating start with the same handle and grant returns the
original attempt and its current outcome. A new attempt requires a new grant.
\repolink{rust/src/olm.rs}{Station lifecycle.}

\begin{samepage}
\paragraph{Attach the peers.}
The supervisor starts one dedicated \path{linkwise-modemd} per peer using the
installed waveform profile. Each child reports ephemeral IPv4 loopback ports
for command, DATA, and virtual PCM in a bounded JSON readiness record. The
supervisor sends a fresh random capability over that child's stdin and
presents it as a prefix on each socket. The child accepts the three
associated sockets within its attachment deadline, closes the listeners,
and serves one wire attempt. The supervisor relays framed 48~kHz mono
samples between the modems. It creates two private temporary mailboxes
and runs the same host-session code used by \path{linkwise mail}. The
listener arms before the caller sends \texttt{CONNECT}. That command begins
the peer handshake inside the existing attempt.
\repolink{rust/src/olm/virtual_bench.rs}{Supervisor and waveform fixture.}
\repolink{rust/src/winlink/host_session.rs}{B2F host session.}\par
\end{samepage}

\paragraph{Observe and verify.}
\path{link_get_status} advances from attachment and peer wait to connected
when both sides report the expected callsigns. The fixture keeps an endpoint's
command and DATA sockets attached when its B2F work finishes so its peer can
observe its own \texttt{DISCONNECTED}. Success requires both B2F client
results, both disconnects, and byte-for-byte agreement between sent and received
fixture envelopes. The record carries the lengths and SHA-256 digests of the
received 23-byte caller attachment and 26-byte responder attachment, plus
the waveform-profile digest. Payload bytes stay off MCP. An eligible RMS
entry establishes a candidate. A virtual modem handshake establishes a
local peer connection. Exact peer mailbox contents establish local B2F
transfer. The gateway is uncontacted in this fixture, and CMS acceptance
and recipient delivery remain unknown. Section~\ref{sec:verification}
describes the associated tests.

\paragraph{End the attempt.}
Command or DATA EOF detaches the pair and ends the affected child. Virtual
audio loss makes the bearer unusable. On abort or deadline, the supervisor
cancels the B2F transports and relay, joins the client work, drops the
temporary mailboxes, and stops and reaps both children. Normal completion
also reaps both children before the station releases its lease. The
connected-abort regression checks cleanup and then verifies an exact
exchange under a fresh grant. If cleanup exceeds its grace period,
the station inhibits further admission while cleanup continues. An
unconfirmed reap or terminal journal failure also inhibits admission.
The shadow release result combines simulated key-down with fixture cleanup.
On restart, journal replay marks an active record interrupted. The
original attempt remains queryable without a recorded release result.
A new shadow process can clear a shadow inhibit because it reconstructs a
simulated radio.
\repolink{rust/tests/olm_station.rs}{Attempt and teardown regressions.}

\section{Commissioned radio design}
\label{sec:live-safety}

\paragraph{One transmit authority.}
The target station verifies actual radio state after preparation and
revalidates it at admission. An authenticated local client or a managed
client reservation can associate with that admission. Each session owns one
radio lease. A session grant authorizes a bounded operation. Each burst
needs a separate expiring, single-use permit checked against the current
lease, verified settings, route, interlocks, and remaining continuous-key
and duty limits. The modem owns protocol timing, and the station supervisor
owns session admission and cancellation. A local arbiter owns physical
key-up through the qualified driver. Independent watchdog release and a
physical emergency stop must still work when the application stalls or
fails. Key-down remains available without a permit.
\repolink{docs/architecture/safety-security-reliability.md}{Safety requirements.}

\stationfigure{Linkwise_Transmit_Boundary}{Proposed transmit authority, separate from protocol timing and independent release controls.}{fig:transmit-boundary}

\begin{samepage}
\paragraph{Control every route.}
The native \path{linkwise-modemd --serve-station} preview can run the modem
waveform engine through CPAL capture and playback. It accepts a legacy local
command and DATA pair outside the shadow coordinator. Physical PTT is
external to the preview, and the setup guide describes radio or interface
VOX and external PTT. The installed MCP service denies key-up and reports
an inactive watchdog. A commissioned design must prevent transmit through
any route outside the arbiter. That includes CAT transmit, serial or
interface keying, audio-triggered VOX, and any modem or client keying
control. Verify the hardware inhibit covers those routes. Loopback binding
limits where clients connect. The live admission service still needs to
verify the client association before allowing use of the radio.
\repolink{rust/docs/modem-setup.md}{Native preview setup.}\par
\end{samepage}

\paragraph{Measure device completion.}
The target session owner gives each playback request a unique transmission
token tied to the session and capture epoch. The route reports acceptance,
presentation, cancellation, and failure. A queue acknowledgment or CPAL
presentation estimate cannot establish when the last sample left the
device. A physical release request and observed PTT down are separate
events. Only calibrated presentation and observed release with the
qualified turnaround allow protocol transmit-complete and the receive
deadline to advance. A late result from a cancelled token is discarded.
An uncertain stop or loss of readback invalidates pending playback,
requests release, and blocks another transmission. Byte accounting
distinguishes pending, peer-acknowledged, discarded, and unresolved bytes.
Unresolved payload is never automatically replayed across sessions.
\repolink{rust/docs/modem-runtime.md}{Current runtime and device contracts.}

\subsection{Fault handling and restart}

The station distinguishes loss of the planning model, loss of the MCP
transport, and loss of a session or device. Each changes a different
owner's authority. In the target design, release proceeds independently
of the application and audit writer.

{\small
\setlength{\LTleft}{0pt}
\setlength{\LTright}{\fill}
\begin{longtable}{@{}>{\raggedright\arraybackslash}p{.21\linewidth}
  >{\raggedright\arraybackslash}p{.34\linewidth}
  >{\raggedright\arraybackslash}p{.38\linewidth}@{}}
\toprule
\textbf{Event and owner} & \textbf{Installed behavior} & \textbf{Commissioned requirement}\\
\midrule
\endfirsthead
\toprule
\textbf{Event and owner} & \textbf{Installed behavior} & \textbf{Commissioned requirement}\\
\midrule
\endhead
\bottomrule
\endfoot
Model or evidence outage.
Station policy owns admission. &
Inventory refresh reports stale or failed input. A running shadow attempt
has its own session deadline. &
The installed live profile specifies whether already authorized work
can finish under continuing safety checks or must abort. Local timers
and release remain available throughout the outage.\\
\addlinespace
MCP stdio EOF.
Station supervisor owns shutdown. &
EOF shuts down the local service and cancels an active station attempt. &
Exercise the local end-of-session policy and measure release when the
control connection ends. Prior authorization continues to bound any
work allowed to finish.\\
\addlinespace
Command, DATA, or audio loss.
Modem and station supervisors own teardown. &
Host EOF detaches the pair. A broken virtual cable ends the managed
exchange and triggers cleanup. &
Invalidate both host attachments and queued output. Confirm presentation
and release state before admitting another client.\\
\addlinespace
Readback mismatch or late completion.
Session owner validates device events. &
Changed simulated settings invalidate preparation. The native preview
uses an estimate of playback presentation. &
Reject stale configuration and playback tokens, cancel pending work,
and hold transmit inhibit while device or release state is uncertain.\\
\addlinespace
Lease or permit expiry, stall, or process restart.
Arbiter and independent watchdog own release. &
Shadow deadlines bound work. Journal replay marks an active attempt
interrupted. &
Force hardware release on the approved deadline. Invalidate outstanding
authority and independently reconcile PTT, radio settings, and client
activity before recovery.\\
\end{longtable}
}

\repolink{rust/src/mcp.rs}{Installed stdio shutdown}.
\repolink{docs/architecture/safety-security-reliability.md}{Target failure handling}.

\paragraph{Commission against the installed route.}
Freeze the application, modem, client, and radio revisions, waveform
digest, audio route, and station profile. Use an inhibited or conducted
setup with independent observation of actual output and PTT state. Measure
presentation, turnaround, abort latency, and watchdog release against
the profile's deadlines. Inject stale readback, lost command or DATA,
device loss, stalled and crashed processes, and a late playback
completion. Record the emergency stop and demonstrate that neither a
fault nor restart can silently re-enable keying. Then run an approved
exchange with an authorized Winlink account, a compatible RMS gateway,
and an agreed test message. The trace must report actual RF and PTT,
expected modem connection, client transfer, CMS acceptance if observed,
and recipient delivery if separately confirmed. Section~\ref{sec:roadmap}
identifies the gates that remain open.
\repolink{docs/architecture/station-integration-status.md}{External qualification gates.}
